\documentclass[11pt,a4paper]{article}
\usepackage[T1]{fontenc}
\usepackage[utf8]{inputenc}
\usepackage{amsmath,amssymb,amsthm}
\usepackage[margin=1in]{geometry}
\usepackage[colorlinks=true,linkcolor=blue,urlcolor=blue]{hyperref}
\theoremstyle{plain}
\newtheorem{theorem}{Theorem}
\newtheorem{lemma}[theorem]{Lemma}
\newtheorem{proposition}[theorem]{Proposition}
\newtheorem{corollary}[theorem]{Corollary}
\theoremstyle{definition}
\newtheorem{remark}[theorem]{Remark}

\title{The empirical recurrence for OEIS A209551, proved by transfer matrix}
\author{Adrian Perez Fontelles\\ \small Independent researcher}
\date{31 August 2026}
\begin{document}
\maketitle

\begin{abstract}
OEIS A209551 counts $(n+1)\times7$ arrays over $\{0,\dots,3\}$ subject to a condition imposed on
every $2\times2$ subblock, and carries an empirical recurrence of order $72$ contributed
by R.~H.~Hardin. It is true, and it is not an empirical matter at all. The condition is local
to each subblock, so the rows of an admissible array are the vertices of a finite digraph
on $S=16384$ vertices whose edges are the admissible consecutive pairs, and the count is a
number of walks: $a(n)=\tfrac1{4}\mathbf{1}^{\!\top}M^{n}\mathbf{1}$. Writing $q$
for the characteristic polynomial of the conjectured recurrence, the recurrence holds for all
$n$ exactly when $\mathbf{1}^{\!\top}M^{j}q(M)\mathbf{1}=0$ for every $j\ge0$, and the
Cayley--Hamilton theorem bounds that to $j<S$. It is therefore a finite computation in exact
integer arithmetic, and it comes out zero.
\end{abstract}

\noindent\small 2020 Mathematics Subject Classification. 05A15, 05B45, 15B36.\normalsize

\section{The sequence and the conjecture}

OEIS A209551 is ``1/4 the number of (n+1)X7 0..3 arrays with every 2X2 subblock having exactly two distinct clockwise edge differences.''. It has offset $1$ and begins
\[
1909, 90186, 3901194, 174401424, 7751834609, 345350598642, 15380031974285, 685056630399243,\ \dots
\]
The entry carries the following comment:
\begin{quote}\small
Empirical: a(n) = 93*a(n-1) -2828*a(n-2) +25792*a(n-3) +371376*a(n-4) -9898959*a(n-5) +50997316*a(n-6) +600219584*a(n-7) -8458938650*a(n-8) +15460753459*a(n-9) +324261513320*a(n-10) -1997445245287*a(n-11) -2888253677269*a(n-12) +61818955362264*a(n-13) -105335609798253*a(n-14) -919549637987857*a(n-15) +3635125259976035*a(n-16) +6249958420274094*a(n-17) -54555499290246451*a(n-18) +11571882298429625*a(n-19) +495953594288336739*a(n-20) -657345762489787062*a(n-21) -2942275874522263687*a(n-22) +6881779713200302370*a(n-23) +11244898913169983370*a(n-24) -42974781730107130109*a(n-25) -23051215421912913536*a(n-26) +185864776798114709968*a(n-27) -12863209686911573111*a(n-28) -587984628431521336548*a(n-29) +276016166528678978015*a(n-30) +1397897524468136739882*a(n-31) -1073664296329817362259*a(n-32) -2537929904791966336740*a(n-33) +2568552806109627886272*a(n-34) +3559706199665273218930*a(n-35) -4354397060651285431293*a(n-36) -3901454743539065189517*a(n-37) +5483143910166447182278*a(n-38) +3391429255331125848034*a(n-39) -5229357377757064837770*a(n-40) -2387498579023401955317*a(n-41) +3802827193424184634130*a(n-42) +1394499277589912306144*a(n-43) -2104672583906607527310*a(n-44) -686223001308051264232*a(n-45) +877657767435453737858*a(n-46) +281766806240431459892*a(n-47) -270362730870675484681*a(n-48) -92957087420024768497*a(n-49) +59494538810619139593*a(n-50) +23369146321510311361*a(n-51) -8826604119797680384*a(n-52) -4236074000925921714*a(n-53) +782901241541277381*a(n-54) +524243697816902673*a(n-55) -26069037492083522*a(n-56) -41764242069671064*a(n-57) -1908566223569902*a(n-58) +1979458374052951*a(n-59) +240266789154528*a(n-60) -47527093603425*a(n-61) -10067754548810*a(n-62) +255750935719*a(n-63) +181717983160*a(n-64) +9192680176*a(n-65) -1053768700*a(n-66) -122161293*a(n-67) -1991984*a(n-68) +258490*a(n-69) +14852*a(n-70) +296*a(n-71) +2*a(n-72)
\end{quote}
As of the ``Last modified'' line on the live entry (Oct 03 2025 15:20:55, revision 9) the statement is
still recorded as empirical, and nothing on the entry records it as settled either way.

Written out, the claim is
\begin{equation}\label{eq:conj}
a(n)\;=\;93\,a(n-1) - 2828\,a(n-2) + 25792\,a(n-3) + 371376\,a(n-4) - 9898959\,a(n-5) + 50997316\,a(n-6) + 600219584\,a(n-7) - 8458938650\,a(n-8) + 15460753459\,a(n-9) + 324261513320\,a(n-10) - 1997445245287\,a(n-11) - 2888253677269\,a(n-12) + 61818955362264\,a(n-13) - 105335609798253\,a(n-14) - 919549637987857\,a(n-15) + 3635125259976035\,a(n-16) + 6249958420274094\,a(n-17) - 54555499290246451\,a(n-18) + 11571882298429625\,a(n-19) + 495953594288336739\,a(n-20) - 657345762489787062\,a(n-21) - 2942275874522263687\,a(n-22) + 6881779713200302370\,a(n-23) + 11244898913169983370\,a(n-24) - 42974781730107130109\,a(n-25) - 23051215421912913536\,a(n-26) + 185864776798114709968\,a(n-27) - 12863209686911573111\,a(n-28) - 587984628431521336548\,a(n-29) + 276016166528678978015\,a(n-30) + 1397897524468136739882\,a(n-31) - 1073664296329817362259\,a(n-32) - 2537929904791966336740\,a(n-33) + 2568552806109627886272\,a(n-34) + 3559706199665273218930\,a(n-35) - 4354397060651285431293\,a(n-36) - 3901454743539065189517\,a(n-37) + 5483143910166447182278\,a(n-38) + 3391429255331125848034\,a(n-39) - 5229357377757064837770\,a(n-40) - 2387498579023401955317\,a(n-41) + 3802827193424184634130\,a(n-42) + 1394499277589912306144\,a(n-43) - 2104672583906607527310\,a(n-44) - 686223001308051264232\,a(n-45) + 877657767435453737858\,a(n-46) + 281766806240431459892\,a(n-47) - 270362730870675484681\,a(n-48) - 92957087420024768497\,a(n-49) + 59494538810619139593\,a(n-50) + 23369146321510311361\,a(n-51) - 8826604119797680384\,a(n-52) - 4236074000925921714\,a(n-53) + 782901241541277381\,a(n-54) + 524243697816902673\,a(n-55) - 26069037492083522\,a(n-56) - 41764242069671064\,a(n-57) - 1908566223569902\,a(n-58) + 1979458374052951\,a(n-59) + 240266789154528\,a(n-60) - 47527093603425\,a(n-61) - 10067754548810\,a(n-62) + 255750935719\,a(n-63) + 181717983160\,a(n-64) + 9192680176\,a(n-65) - 1053768700\,a(n-66) - 122161293\,a(n-67) - 1991984\,a(n-68) + 258490\,a(n-69) + 14852\,a(n-70) + 296\,a(n-71) + 2\,a(n-72)
\end{equation}
for all $n>72$.

\section{The condition is local, so the rows form a digraph}

Write an array of the shape in question as its list of rows
$r^{(1)},r^{(2)},\dots$, each an element of
$V=\{0,1,\dots,3\}^{7}$, so $|V|=S=16384$. A $2\times2$ subblock of the array
occupies two consecutive rows and two consecutive columns; if the two rows are $r$
and $s$ (in that order) and the two columns are numbered $j,j+1$,
the subblock is
\[
\begin{pmatrix}\alpha&\beta\\\gamma&\delta\end{pmatrix}=\begin{pmatrix}r_j&r_{j+1}\\s_j&s_{j+1}\end{pmatrix}.
\]
The entry's requirement on that subblock is
\begin{equation}\label{eq:loc}
\#\{\beta-\alpha,\ \delta-\beta,\ \gamma-\delta,\ \alpha-\gamma\}\in\{2\},
\end{equation}
a condition on the four entries alone.

For $r,s\in V$ declare $r\to s$ an edge of a digraph $G$ when, for every
$1\le j\le 7-1$,
\begin{itemize}
\item the four entries above satisfy \eqref{eq:loc};

\end{itemize}
Let $M\in\{0,1\}^{S\times S}$ be the adjacency matrix of $G$ and $\mathbf{1}$ the
all-ones vector.

\begin{lemma}\label{lem:walk}
The arrays counted by the entry with $L$ rows are exactly the walks
$r^{(1)}\to r^{(2)}\to\cdots\to r^{(L)}$ of length $L-1$ in $G$. Consequently
\[
a(n)\;=\;\tfrac1{4}\mathbf{1}^{\!\top}M^{n}\mathbf{1} .
\]
\end{lemma}

\begin{proof}
An array with $L$ rows and $7$ columns is precisely a sequence
$r^{(1)},\dots,r^{(L)}$ of elements of $V$; conversely any such sequence is an array.
Every $2\times2$ subblock lies in two consecutive rows $r^{(t)},r^{(t+1)}$ and two
consecutive columns, so the array satisfies the entry's condition on all of its subblocks
if and only if each consecutive pair $\bigl(r^{(t)},r^{(t+1)}\bigr)$ satisfies it for
every column pair --- which is exactly the edge relation of $G$. The number of walks of
length $L-1$, summed over all starting and ending vertices, is
$\mathbf{1}^{\!\top}M^{L-1}\mathbf{1}$, since $(M^{k})_{r,s}$ counts the walks of
length $k$ from $r$ to $s$. The entry's index $n$ corresponds to $L=n+1$ rows. The entry counts $\tfrac1{4}$ of those arrays, a constant factor that passes through every step below unchanged.
\end{proof}

\section{The criterion}

Let $q(t)=t^{72}-\sum_i c_i\,t^{72-i}$ be the characteristic polynomial of
\eqref{eq:conj}, with the $c_i$ read off that equation.

\begin{lemma}\label{lem:crit}
For every $n$ with $n+1-1\ge72$,
\[
a(n)-\sum_i c_i\,a(n-i)\;=\;\tfrac1{4}\mathbf{1}^{\!\top}M^{\,n+1-1-72}\,q(M)\,\mathbf{1} .
\]
Hence \eqref{eq:conj} holds from that point on if and only if
$\mathbf{1}^{\!\top}M^{j}q(M)\mathbf{1}=0$ for every $j\ge0$; and it is enough to check
$j=0,1,\dots,S-1$.
\end{lemma}

\begin{proof}
By Lemma~\ref{lem:walk}, $a(n-i)=\tfrac1{4}\mathbf{1}^{\!\top}M^{\,n+1-1-i}\mathbf{1}$, so
\[
a(n)-\sum_i c_i a(n-i)
=\tfrac1{4}\mathbf{1}^{\!\top}\Bigl(M^{m}-\sum_i c_i M^{m-i}\Bigr)\mathbf{1}
=\tfrac1{4}\mathbf{1}^{\!\top}M^{\,m-72}\,q(M)\,\mathbf{1}
\]
with $m=n+1-1$, which is the stated identity; the scalar $\tfrac1{4}$ never vanishes, so
it does not affect whether the left side is zero. The equivalence follows by letting
$m-72=j$ run over $\ge0$. For the last clause put $w=q(M)\mathbf{1}$. By the
Cayley--Hamilton theorem $M^{S}$ is an integer combination of $I,M,\dots,M^{S-1}$, so
$\mathbf{1}^{\!\top}M^{j}w$ for $j\ge S$ is the corresponding combination of the earlier
values; if those all vanish, so do all the rest.
\end{proof}

\section{The computation}

\begin{theorem}
The recurrence \eqref{eq:conj} holds for every $n>72$, which is the statement of
Section~1.
\end{theorem}

\begin{proof}
The digraph was constructed from \eqref{eq:loc} by a depth-first scan across the $7$
columns: the edge condition is a conjunction of one constraint per consecutive column
pair, so the successors of a row $r$ are enumerated column by column without testing all
$S^2$ pairs. The vector $w=q(M)\mathbf{1}\in\mathbb{Z}^{S}$ was then formed by $72$
matrix--vector products in exact integer arithmetic, and $\mathbf{1}^{\!\top}M^{j}w$ was
evaluated for $j=0,1,\dots$, again exactly. Every one of those integers vanishes from the
index corresponding to $n=72+1$ onwards, and the run continued until $w$ itself was the
zero vector, which settles every larger $j$ at once. By Lemma~\ref{lem:crit} the recurrence
holds for every $n>72$.
\end{proof}

\begin{remark}
The argument gives more than the recurrence: it shows the sequence is $C$-finite with
characteristic polynomial dividing that of $M$, so it has a rational generating function
whose denominator divides $\det(I-xM)$. The conjectured recurrence is one consequence; the
minimal one is a divisor of $q$.
\end{remark}

\section{Verification}

Three checks, each independent of the algebra above.

First, the digraph was built directly from the entry's stated condition and the walk counts
$\tfrac1{4}\mathbf{1}^{\!\top}M^{L-1}\mathbf{1}$ compared against the entry's DATA at
every one of the $13$ published indices; the values agree exactly. That is what ties
the digraph to the sequence: a model that reproduced only some of the published terms would
be the wrong model, and would have been rejected. In particular it is what rules out a
misreading of the entry's wording, since a different reading of the condition gives a
different digraph and different counts.

Second, the recurrence \eqref{eq:conj} was evaluated on those published terms in exact
integer arithmetic, with no matrices involved, and vanishes wherever it applies.

Third, the annihilation test of Lemma~\ref{lem:crit} was run in exact integer arithmetic
until the working vector was identically zero, not to a sampled prefix, and returned zero
throughout.

\begin{thebibliography}{9}
\bibitem{oeis} The OEIS Foundation, \emph{The On-Line Encyclopedia of Integer
Sequences}, \url{https://oeis.org}, 2026. Sequence A209551.
\bibitem{stanley} R.~P.~Stanley, \emph{Enumerative Combinatorics}, Volume 1, 2nd ed.,
Cambridge University Press, 2011. (Transfer-matrix method, Section 4.7.)
\bibitem{fs} P.~Flajolet and R.~Sedgewick, \emph{Analytic Combinatorics}, Cambridge
University Press, 2009.
\bibitem{horn} R.~A.~Horn and C.~R.~Johnson, \emph{Matrix Analysis}, 2nd ed., Cambridge
University Press, 2013.
\end{thebibliography}

\end{document}
